\documentclass[10pt,a4paper]{amsart}
\usepackage{amsmath,amssymb,amsthm,graphicx}
% Retained preamble for the Stacks Project chapter divisors.tex (Divisors).
% The chapter file pinned-original-source.tex carries only \input{preamble} on its first line, so no
% preamble text lives inside the chapter itself.  The Stacks Project uses one single shared preamble for
% every chapter; the snapshot stages it verbatim in sources/stacks-chow/preamble.tex.  The block below is
% copied mechanically, line for line, from that file's lines 14-18 (the Xy-pic setup) and lines 154-217
% (the native theorem declarations together with the mathematical macros).  Nothing was added, removed or
% redefined.  The wrapper class and the amsmath/amssymb/amsthm/graphicx base packages are supplied by the
% assembler, so the class/\externaldocument part of the shared preamble is not repeated here.
% For commutative diagrams we use Xy-pic
\usepackage[all]{xy}

% We use 2cell for 2-commutative diagrams.
\xyoption{2cell}
% Theorem environments.
%
\theoremstyle{plain}
\newtheorem{theorem}[subsection]{Theorem}
\newtheorem{proposition}[subsection]{Proposition}
\newtheorem{lemma}[subsection]{Lemma}

\theoremstyle{definition}
\newtheorem{definition}[subsection]{Definition}
\newtheorem{example}[subsection]{Example}
\newtheorem{exercise}[subsection]{Exercise}
\newtheorem{situation}[subsection]{Situation}

\theoremstyle{remark}
\newtheorem{remark}[subsection]{Remark}
\newtheorem{remarks}[subsection]{Remarks}

\numberwithin{equation}{subsection}

% Macros
%
\def\lim{\mathop{\mathrm{lim}}\nolimits}
\def\colim{\mathop{\mathrm{colim}}\nolimits}
\def\Spec{\mathop{\mathrm{Spec}}}
\def\Hom{\mathop{\mathrm{Hom}}\nolimits}
\def\Ext{\mathop{\mathrm{Ext}}\nolimits}
\def\SheafHom{\mathop{\mathcal{H}\!\mathit{om}}\nolimits}
\def\SheafExt{\mathop{\mathcal{E}\!\mathit{xt}}\nolimits}
\def\Sch{\mathit{Sch}}
\def\Mor{\mathop{\mathrm{Mor}}\nolimits}
\def\Ob{\mathop{\mathrm{Ob}}\nolimits}
\def\Sh{\mathop{\mathit{Sh}}\nolimits}
\def\NL{\mathop{N\!L}\nolimits}
\def\CH{\mathop{\mathrm{CH}}\nolimits}
\def\proetale{{pro\text{-}\acute{e}tale}}
\def\etale{{\acute{e}tale}}
\def\QCoh{\mathit{QCoh}}
\def\Ker{\mathop{\mathrm{Ker}}}
\def\Im{\mathop{\mathrm{Im}}}
\def\Coker{\mathop{\mathrm{Coker}}}
\def\Coim{\mathop{\mathrm{Coim}}}

% Boxtimes
%
\DeclareMathSymbol{\boxtimes}{\mathbin}{AMSa}{"02}

%
% Macros for moduli stacks/spaces
%
\def\QCohstack{\mathcal{QC}\!\mathit{oh}}
\def\Cohstack{\mathcal{C}\!\mathit{oh}}
\def\Spacesstack{\mathcal{S}\!\mathit{paces}}
\def\Quotfunctor{\mathrm{Quot}}
\def\Hilbfunctor{\mathrm{Hilb}}
\def\Curvesstack{\mathcal{C}\!\mathit{urves}}
\def\Polarizedstack{\mathcal{P}\!\mathit{olarized}}
\def\Complexesstack{\mathcal{C}\!\mathit{omplexes}}
% \Pic is the operator that assigns to X its picard group, usage \Pic(X)
% \Picardstack_{X/B} denotes the Picard stack of X over B
% \Picardfunctor_{X/B} denotes the Picard functor of X over B
\def\Pic{\mathop{\mathrm{Pic}}\nolimits}
\def\Picardstack{\mathcal{P}\!\mathit{ic}}
\def\Picardfunctor{\mathrm{Pic}}
\def\Deformationcategory{\mathcal{D}\!\mathit{ef}}

% --- appended 2026-10-12 (Chief): the project own definitions of the reference/slogan/history
% display environments, copied mechanically from the shared Stacks preamble lines 8-12
% (\usepackage{verbatim} plus the three \newenvironment lines, which the chapter preamble keeps
% verbatim from the project preamble). Nothing added, removed or redefined: these environments
% carry no mathematical content and the project itself typesets them as comments. Without these
% three lines a counted statement whose own text uses them fails with Environment slogan undefined.
\usepackage{verbatim}
\newenvironment{reference}{\comment}{\endcomment}
\newenvironment{slogan}{\comment}{\endcomment}
\newenvironment{history}{\comment}{\endcomment}
\usepackage{cleveref}
\crefname{originallemma}{original lemma}{original lemmas}
\Crefname{originallemma}{Original lemma}{Original lemmas}
\makeatletter
\newlabel{algebra-lemma-flat-finite-presentation-limit-flat}{{02JO}{source}{}{sourcepointer-1}{}}
\newlabel{algebra-lemma-flat-finite-presentation-limit-flat@cref}{{[originallemma][0][]02JO}{[1][0][]source}{}{}{}}
\makeatother
\crefname{originallemma}{original lemma}{original lemmas}
\Crefname{originallemma}{Original lemma}{Original lemmas}
\makeatletter
\newlabel{algebra-lemma-grothendieck}{{00MF}{source}{}{sourcepointer-2}{}}
\newlabel{algebra-lemma-grothendieck@cref}{{[originallemma][0][]00MF}{[1][0][]source}{}{}{}}
\makeatother
\crefname{originallemma}{original lemma}{original lemmas}
\Crefname{originallemma}{Original lemma}{Original lemmas}
\makeatletter
\newlabel{algebra-lemma-local-isomorphism}{{00QR}{source}{}{sourcepointer-3}{}}
\newlabel{algebra-lemma-local-isomorphism@cref}{{[originallemma][0][]00QR}{[1][0][]source}{}{}{}}
\makeatother
\crefname{originallemma}{original lemma}{original lemmas}
\Crefname{originallemma}{Original lemma}{Original lemmas}
\makeatletter
\newlabel{algebra-lemma-mod-injective-general}{{046Y}{source}{}{sourcepointer-4}{}}
\newlabel{algebra-lemma-mod-injective-general@cref}{{[originallemma][0][]046Y}{[1][0][]source}{}{}{}}
\makeatother
\crefname{originallemma}{original lemma}{original lemmas}
\Crefname{originallemma}{Original lemma}{Original lemmas}
\makeatletter
\newlabel{algebra-lemma-regular-sequence-in-neighbourhood}{{061L}{source}{}{sourcepointer-5}{}}
\newlabel{algebra-lemma-regular-sequence-in-neighbourhood@cref}{{[originallemma][0][]061L}{[1][0][]source}{}{}{}}
\makeatother
\crefname{originaltheorem}{original theorem}{original theorems}
\Crefname{originaltheorem}{Original theorem}{Original theorems}
\makeatletter
\newlabel{algebra-theorem-openness-flatness}{{00RC}{source}{}{sourcepointer-6}{}}
\newlabel{algebra-theorem-openness-flatness@cref}{{[originaltheorem][0][]00RC}{[1][0][]source}{}{}{}}
\makeatother
\crefname{originallemma}{original lemma}{original lemmas}
\Crefname{originallemma}{Original lemma}{Original lemmas}
\makeatletter
\newlabel{lemma-relative-Cartier}{{056Q}{source}{}{sourcepointer-7}{}}
\newlabel{lemma-relative-Cartier@cref}{{[originallemma][0][]056Q}{[1][0][]source}{}{}{}}
\makeatother
\crefname{originallemma}{original lemma}{original lemmas}
\Crefname{originallemma}{Original lemma}{Original lemmas}
\makeatletter
\newlabel{morphisms-lemma-composition-finite-presentation}{{01TR}{source}{}{sourcepointer-8}{}}
\newlabel{morphisms-lemma-composition-finite-presentation@cref}{{[originallemma][0][]01TR}{[1][0][]source}{}{}{}}
\makeatother
\crefname{originallemma}{original lemma}{original lemmas}
\Crefname{originallemma}{Original lemma}{Original lemmas}
\makeatletter
\newlabel{morphisms-lemma-finite-presentation-permanence}{{02FV}{source}{}{sourcepointer-9}{}}
\newlabel{morphisms-lemma-finite-presentation-permanence@cref}{{[originallemma][0][]02FV}{[1][0][]source}{}{}{}}
\makeatother
\title{A closed subscheme whose fibres are Cartier divisors is locally contained in a relative effective Cartier divisor, and agrees with one under flatness or local principality}
\author{The Stacks Project Authors}
\date{stacks-project:divisors.tex:local-snapshot-20261012; GNU FDL 1.2 or later}
\begin{document}
\maketitle
\noindent\textit{Editorial note.} All delivered mathematical text is the author's own, copied from the pinned Stacks Project chapter file pinned-original-source.tex; no statement or proof step was rewritten, reordered or removed.\par\medskip
\noindent\textit{Editorial note.} Following the multi-claim convention for this source, source\_url is the source-level URL https://stacks.math.columbia.edu/ and the source label lemma-fibre-Cartier is carried in claim\_key and here rather than in a per-claim tag URL.\par\medskip
\noindent\textit{Editorial note.} The preamble is the chapter's shared Stacks Project preamble, retained as sources/stacks-divisors/retained-preamble.tex: a line-for-line excerpt of the shared Stacks preamble staged with this snapshot. No macro was invented or redefined, and this card adds no preamble text of its own.\par\medskip
\noindent\textit{Editorial note.} Statement span 4190-4213 and proof span 4215-4291 come from the halfA pool card operations/pool-cards/stacks-divisors-halfA.json and were re-read in the pinned chapter file before this card was written.\par\medskip
\noindent\textit{Editorial note.} The proof contains a parenthetical mismatch of parentheses in the base change display of the fibre ring and writes 'essentially of finite presentation' in its own wording; both are the author's text and are preserved unchanged.\par\medskip
\noindent\textit{Editorial note.} 9 cross-references occur inside the counted spans; they are listed in external\_references as navigation pointers only.\par\medskip
\noindent\textit{Allowed standard prerequisites (editorial).}\begin{itemize}
\item Relative effective Cartier divisors and their stability under base change, cited from the same chapter. Cited at proof lines 4258-4259. Transfers the construction performed on an approximating Noetherian model back to the original situation.
\item Approximation of a flat and finitely presented algebra by a flat finite type algebra over a Noetherian ring, Grothendieck's generic freeness and nonzerodivisor statement, the existence of a neighbourhood on which a regular element stays regular with flat quotient, and the openness of flatness, all cited from the algebra chapter. Cited at proof lines 4234-4252. These four inputs produce the element g outside the prime about which the constructed divisor is a relative effective Cartier divisor.
\item A surjective map of flat local algebras essentially of finite presentation that becomes an isomorphism after dividing by the maximal ideal is an isomorphism, and the local criterion for an isomorphism in a neighbourhood. Cited at proof lines 4267-4272. Upgrades the containment of part (1) to an equality under the flatness hypothesis of part (2).
\item The equivalence between the source morphism of a closed subscheme being locally of finite presentation and the composite being so, cited from the morphisms chapter. Cited at proof lines 4288-4290. Deduces the final statements of the lemma from parts (2) and (3).
\end{itemize}\medskip
\ifdefined\hypertarget\hypertarget{sourcepointer-1}{}\fi
\noindent\textit{Source reference.} Original lemma 02JO: focused-seeds/AGP/sources/stacks-snapshot-20261011/tags.txt, source lines 3343--3343. Same-chapter reference used in the retained author text at source line(s) 4106, 4234, 5430. It is supplied only as a navigation pointer so that the retained text compiles; the counted item delivers the author's own statement and the author's complete proof, and no part of the referenced item is copied in or substituted for the counted argument.\par\medskip
\ifdefined\hypertarget\hypertarget{sourcepointer-2}{}\fi
\noindent\textit{Source reference.} Original lemma 00MF: focused-seeds/AGP/sources/stacks-snapshot-20261011/tags.txt, source lines 821--821. Same-chapter reference used in the retained author text at source line(s) 4246. It is supplied only as a navigation pointer so that the retained text compiles; the counted item delivers the author's own statement and the author's complete proof, and no part of the referenced item is copied in or substituted for the counted argument.\par\medskip
\ifdefined\hypertarget\hypertarget{sourcepointer-3}{}\fi
\noindent\textit{Source reference.} Original lemma 00QR: focused-seeds/AGP/sources/stacks-snapshot-20261011/tags.txt, source lines 977--977. Same-chapter reference used in the retained author text at source line(s) 4274, 5560. It is supplied only as a navigation pointer so that the retained text compiles; the counted item delivers the author's own statement and the author's complete proof, and no part of the referenced item is copied in or substituted for the counted argument.\par\medskip
\ifdefined\hypertarget\hypertarget{sourcepointer-4}{}\fi
\noindent\textit{Source reference.} Original lemma 046Y: focused-seeds/AGP/sources/stacks-snapshot-20261011/tags.txt, source lines 5489--5489. Same-chapter reference used in the retained author text at source line(s) 4271, 5557. It is supplied only as a navigation pointer so that the retained text compiles; the counted item delivers the author's own statement and the author's complete proof, and no part of the referenced item is copied in or substituted for the counted argument.\par\medskip
\ifdefined\hypertarget\hypertarget{sourcepointer-5}{}\fi
\noindent\textit{Source reference.} Original lemma 061L: focused-seeds/AGP/sources/stacks-snapshot-20261011/tags.txt, source lines 7765--7765. Same-chapter reference used in the retained author text at source line(s) 2970, 4250, 4869, 5359, 5534, 5635. It is supplied only as a navigation pointer so that the retained text compiles; the counted item delivers the author's own statement and the author's complete proof, and no part of the referenced item is copied in or substituted for the counted argument.\par\medskip
\ifdefined\hypertarget\hypertarget{sourcepointer-6}{}\fi
\noindent\textit{Source reference.} Original theorem 00RC: focused-seeds/AGP/sources/stacks-snapshot-20261011/tags.txt, source lines 998--998. Same-chapter reference used in the retained author text at source line(s) 4132, 4252, 5481. It is supplied only as a navigation pointer so that the retained text compiles; the counted item delivers the author's own statement and the author's complete proof, and no part of the referenced item is copied in or substituted for the counted argument.\par\medskip
\ifdefined\hypertarget\hypertarget{sourcepointer-7}{}\fi
\noindent\textit{Source reference.} Original lemma 056Q: focused-seeds/AGP/sources/stacks-snapshot-20261011/tags.txt, source lines 6718--6718. Same-chapter reference used in the retained author text at source line(s) 4259, 5369. It is supplied only as a navigation pointer so that the retained text compiles; the counted item delivers the author's own statement and the author's complete proof, and no part of the referenced item is copied in or substituted for the counted argument.\par\medskip
\ifdefined\hypertarget\hypertarget{sourcepointer-8}{}\fi
\noindent\textit{Source reference.} Original lemma 01TR: focused-seeds/AGP/sources/stacks-snapshot-20261011/tags.txt, source lines 2384--2384. Same-chapter reference used in the retained author text at source line(s) 4289, 5592. It is supplied only as a navigation pointer so that the retained text compiles; the counted item delivers the author's own statement and the author's complete proof, and no part of the referenced item is copied in or substituted for the counted argument.\par\medskip
\ifdefined\hypertarget\hypertarget{sourcepointer-9}{}\fi
\noindent\textit{Source reference.} Original lemma 02FV: focused-seeds/AGP/sources/stacks-snapshot-20261011/tags.txt, source lines 3206--3206. Same-chapter reference used in the retained author text at source line(s) 4290, 5593. It is supplied only as a navigation pointer so that the retained text compiles; the counted item delivers the author's own statement and the author's complete proof, and no part of the referenced item is copied in or substituted for the counted argument.\par\medskip
\begin{lemma}
\label{lemma-fibre-Cartier}
Let $\varphi : X \to S$ be a flat morphism which is locally of finite
presentation. Let $Z \subset X$ be a closed subscheme.
Let $x \in Z$ with image $s \in S$.
\begin{enumerate}
\item If $Z_s \subset X_s$ is a Cartier divisor in a neighbourhood of $x$,
then there exists an open $U \subset X$ and a
relative effective Cartier divisor $D \subset U$ such that
$Z \cap U \subset D$ and $Z_s \cap U = D_s$.
\item If $Z_s \subset X_s$ is a Cartier divisor in a neighbourhood of $x$,
the morphism $Z \to X$ is of finite presentation, and $Z \to S$ is flat at
$x$, then we can choose $U$ and $D$ such that $Z \cap U = D$.
\item If $Z_s \subset X_s$ is a Cartier divisor in a neighbourhood of $x$
and $Z$ is a locally principal closed subscheme of $X$ in a neighbourhood
of $x$, then we can choose $U$ and $D$ such that $Z \cap U = D$.
\end{enumerate}
In particular, if $Z \to S$ is locally of finite presentation and flat and
all fibres $Z_s \subset X_s$ are effective Cartier divisors, then
$Z$ is a relative effective Cartier divisor. Similarly, if $Z$
is a locally principal closed subscheme of $X$ such that all fibres
$Z_s \subset X_s$ are effective Cartier divisors, then
$Z$ is a relative effective Cartier divisor.
\end{lemma}

\begin{proof}
Choose affine open neighbourhoods $\Spec(A)$ of $s$ and
$\Spec(B)$ of $x$ such that
$\varphi(\Spec(B)) \subset \Spec(A)$.
Let $\mathfrak p \subset A$ be the prime ideal corresponding to $s$.
Let $\mathfrak q \subset B$ be the prime ideal corresponding to $x$.
Let $I \subset B$ be the ideal corresponding to $Z$.
By the initial assumption of the lemma we know that
$A \to B$ is flat and of finite presentation.
The assumption in (1) means that, after shrinking $\Spec(B)$, we may
assume $I(B \otimes_A \kappa(\mathfrak p))$ is generated by a single
element which is a nonzerodivisor in $B \otimes_A \kappa(\mathfrak p)$.
Say $f \in I$ maps to this generator. We claim that after inverting
an element $g \in B$, $g \not \in \mathfrak q$ the closed subscheme
$D = V(f) \subset \Spec(B_g)$ is a relative effective Cartier
divisor.

\medskip\noindent
By
Algebra, Lemma \ref{algebra-lemma-flat-finite-presentation-limit-flat}
we can find a flat finite type ring map $A_0 \to B_0$ of Noetherian
rings, an element $f_0 \in B_0$, a ring map $A_0 \to A$ and an isomorphism
$A \otimes_{A_0} B_0 \cong B$. If $\mathfrak p_0 = A_0 \cap \mathfrak p$
then we see that
$$
B \otimes_A \kappa(\mathfrak p) =
\left(B_0 \otimes_{A_0} \kappa(\mathfrak p_0)\right)
\otimes_{\kappa(\mathfrak p_0))} \kappa(\mathfrak p)
$$
hence $f_0$ is a nonzerodivisor in $B_0 \otimes_{A_0} \kappa(\mathfrak p_0)$.
By
Algebra, Lemma \ref{algebra-lemma-grothendieck}
we see that $f_0$ is a nonzerodivisor in $(B_0)_{\mathfrak q_0}$
where $\mathfrak q_0 = B_0 \cap \mathfrak q$ and
that $(B_0/f_0B_0)_{\mathfrak q_0}$ is flat over $A_0$. Hence by
Algebra, Lemma \ref{algebra-lemma-regular-sequence-in-neighbourhood}
and
Algebra, Theorem \ref{algebra-theorem-openness-flatness}
there exists a $g_0 \in B_0$, $g_0 \not \in \mathfrak q_0$ such
that $f_0$ is a nonzerodivisor in $(B_0)_{g_0}$ and such that
$(B_0/f_0B_0)_{g_0}$ is flat over $A_0$. Hence we see that
$D_0 = V(f_0) \subset \Spec((B_0)_{g_0})$ is a relative effective
Cartier divisor. Since we know that this property is preserved under
base change, see
Lemma \ref{lemma-relative-Cartier},
we obtain the claim mentioned above with $g$ equal to the image of $g_0$
in $B$.

\medskip\noindent
At this point we have proved (1). To see (2) consider the closed
immersion $Z \to D$. The surjective ring map
$u : \mathcal{O}_{D, x} \to \mathcal{O}_{Z, x}$
is a map of flat local $\mathcal{O}_{S, s}$-algebras which
are essentially of finite presentation, and which becomes an
isomorphisms after dividing by $\mathfrak m_s$. Hence it is
an isomorphism, see
Algebra, Lemma \ref{algebra-lemma-mod-injective-general}.
It follows that $Z \to D$ is an isomorphism in a neighbourhood
of $x$, see
Algebra, Lemma \ref{algebra-lemma-local-isomorphism}.
To see (3), after possibly shrinking $U$ we may assume that
the ideal of $D$ is generated by a single nonzerodivisor $f$
and the ideal of $Z$ is generated by an element $g$. Then
$f = gh$. But $g|_{U_s}$ and $f|_{U_s}$ cut out the same
effective Cartier divisor in a neighbourhood of $x$. Hence
$h|_{X_s}$ is a unit in $\mathcal{O}_{X_s, x}$, hence $h$ is
a unit in $\mathcal{O}_{X, x}$ hence $h$ is a unit in an
open neighbourhood of $x$. I.e., $Z \cap U = D$ after shrinking $U$.

\medskip\noindent
The final statements of the lemma follow immediately from
parts (2) and (3), combined with the fact that $Z \to S$
is locally of finite presentation if and only if $Z \to X$ is
of finite presentation, see
Morphisms, Lemmas \ref{morphisms-lemma-composition-finite-presentation} and
\ref{morphisms-lemma-finite-presentation-permanence}.
\end{proof}

\end{document}
